% =============================================================
% Section: Validation
% =============================================================

\section{Calibration checks and directional benchmarks}\label{sec:validation}

The checks in this section have a narrow role. The model matches the calibration targets by construction. The external benchmarks are directional plausibility checks, not causal validation of the counterfactual. Magnitudes from the Portuguese 1996 reform and the Brazilian 1988 reform should be compared in sign and order of magnitude, not in level.

Panel B1 evaluates the model against domestic moments that are not used to fit the main parameters. The model is close on mean formal hours and on the share of formal workers above $36$h. It understates the cross-sector range in average formal hours, which is a limitation of the parsimonious national hours distribution. As a check, the replication package perturbs the formal hours distribution $\theta$ within a $\pm 5$ pp envelope around the DIEESE 2024 weights and re-runs the calibration. The fitted Panel B1 moments shift by $0.4$ to $0.6$ h/week and $1$ to $2$ pp, with the sign of the gap preserved. These non-targeted moments are informative, but they are not a separate identification exercise.

Panel B2 compares the model with two hours reforms. The Portuguese 1996 reform has a common $44\to 40$h cap, and the Brazilian 1988 reform has a $48\to 44$h cap. Signs match. Aggregate output falls and hourly productivity rises in both data and model. Magnitudes do not match. The firm-level Portuguese estimate of $+4.4\%$ in hourly productivity is about $2.6$ points above the aggregate model prediction, a gap consistent with worker selection, work intensification, and sectoral capital deepening that the model does not encode. The Brazilian 1988 evidence points in the same direction, with a slightly larger empirical hourly-wage response. These comparisons are directional benchmarks rather than lower bounds for Brazil's effective $A_{\text{req}}$. Portugal and Brazil 1988 differ from the current Brazilian 36h/40h policy environment in informality, institutions, and reform magnitude.

\begin{table}[htbp]
\centering
\caption{Calibration checks and directional benchmarks.}\label{tab:calibration_fit}
\small
\begin{tabular}{lccc}
\hline\hline
Moment & Observed & Model & Gap \\
\hline
\multicolumn{4}{l}{\textit{Panel A: Calibration targets (matched by construction)}} \\
\hline
Informality rate, small firms ($\leq 49$ empl.)  & 0.500 & 0.500 & 0.000 \\
Informality rate, large firms                    & 0.200 & 0.200 & 0.000 \\
Formal employment share, small firms             & 0.590 & 0.590 & 0.000 \\
Formal employment share, large firms             & 0.410 & 0.410 & 0.000 \\
Aggregate informality, narrow PNAD               & 0.378 & 0.377 & 0.001 \\
Wage premium $R$, MNR (2015) central target       & 1.400 & 1.400 & 0.000 \\
\hline
\multicolumn{4}{l}{\textit{Panel B1: Domestic checks outside calibration}} \\
\hline
Average weekly formal hours, habitual measure            & 41.23 & 42.24 & 1.01 h/week \\
Formal share above $36$h, habitual measure (\%)          & 85.7  & 91.5  & 5.80 pp \\
Sectoral range in average formal hours                   & 3.31  & 0.80  & $-2.51$ h/week \\
\hline
\multicolumn{4}{l}{\textit{Panel B2: External directional benchmarks}} \\
\hline
$\Delta Y$ at $44\to 40$h cap, Portugal 1996 (firm-level) & $-3.20$ & $-2.00$ & 1.20 pp \\
$\Delta Y/h$ at $44\to 40$h cap, Portugal 1996 (\%)       & $+4.40$ & $+1.77$ & 2.63 pp \\
$\Delta W/h$ at $48\to 44$h, Brazil 1988 (\%, GMC 2003)   & $+7.00$ & $+1.77$ & 5.23 pp \\
\hline\hline
\end{tabular}
\\[2pt]
\begin{minipage}{0.95\textwidth}
\justifying
{\footnotesize Panel B1 rows 1 and 2 compare PNAD habitual-hours moments with the baseline model moments implied by the contracted-hours calibration. Row 3 reports the range across agriculture, industry, and services in average formal weekly hours. The model value comes from the sectoral extension using the same hours-bin structure, so the negative gap means that the model understates observed sectoral dispersion in formal hours. The aggregate-hours gap is driven mainly by the maintained assumption $h_I=44$ for informal workers and is treated as a sensitivity margin in the online appendix, not as a validation target ($A_{\text{req}}$ shifts by $<0.2$ pp for $h_I\in[38,46]$). Panel B2 compares the model's aggregate prediction at the $44\to 40$h cap to the Portuguese firm-level estimate \citep{AsaiLopesTondini2024} and to Brazil's 1988 reform of comparable cap-cut magnitude (48$\to$44h, $\sim$8\%) as documented by \citet{GonzagaMenezesFilhoCamargo2003}. Both external benchmarks suggest the model is conservative on the hourly-productivity response. The model's $e(h)$ captures fatigue and not worker selection, work intensification, or sectoral capital deepening. Sources: PNAD/SIDRA, RAIS 2022, DIEESE 2024, \citet{MeghirNaritaRobin2015}, \citet{AsaiLopesTondini2024}, and \citet{GonzagaMenezesFilhoCamargo2003} \citep{PNAD2025,RAIS2022,DIEESE2024}.}
\end{minipage}
\end{table}
